\documentclass[aspectratio=169,11pt]{beamer}

\usepackage[T1]{fontenc}
\usepackage{lmodern}
\usepackage{amsmath,amssymb}
\usepackage{booktabs}
\usepackage{graphicx}

\definecolor{Ink}{HTML}{16283A}
\definecolor{Teal}{HTML}{087F8C}
\definecolor{Coral}{HTML}{D95D39}
\definecolor{Paper}{HTML}{F7F4EE}
\definecolor{Muted}{HTML}{526475}
\definecolor{PaleTeal}{HTML}{DCEBED}
\definecolor{PaleCoral}{HTML}{F3E1DA}

\setbeamercolor{normal text}{fg=Ink,bg=Paper}
\setbeamercolor{frametitle}{fg=Ink,bg=Paper}
\setbeamercolor{structure}{fg=Teal}
\setbeamercolor{alerted text}{fg=Coral}
\setbeamercolor{title}{fg=Paper}
\setbeamercolor{subtitle}{fg=Paper}
\setbeamercolor{author}{fg=Paper}
\setbeamercolor{date}{fg=Paper}
\setbeamerfont{title}{series=\bfseries,size=\fontsize{26}{30}}
\setbeamerfont{subtitle}{size=\normalsize}
\setbeamerfont{frametitle}{series=\bfseries,size=\fontsize{18}{21}}
\setbeamertemplate{navigation symbols}{}
\setbeamersize{text margin left=0.65cm,text margin right=0.65cm}

\setbeamertemplate{frametitle}{
  \vspace{0.27cm}
  \insertframetitle\par
  \vspace{0.06cm}
  {\color{Teal}\rule{\textwidth}{0.8pt}}
}

\setbeamertemplate{footline}{
  \leavevmode
  \hbox{%
    \begin{beamercolorbox}[wd=.88\paperwidth,ht=2.8ex,dp=1.2ex,leftskip=.65cm]{normal text}
      \scriptsize\color{Muted}\texttt{github.com/amchavezu/ai-05-ide}
    \end{beamercolorbox}%
    \begin{beamercolorbox}[wd=.12\paperwidth,ht=2.8ex,dp=1.2ex,center]{normal text}
      \scriptsize\color{Muted}\insertframenumber/14
    \end{beamercolorbox}%
  }
}

\newcommand{\key}[1]{\textcolor{Teal}{\textbf{#1}}}
\newcommand{\caveat}[1]{\textcolor{Coral}{\textbf{#1}}}
\newcommand{\sourcefoot}{\scriptsize\color{Muted}Source: Ide and Talam\`as (2025), accepted manuscript v12.}

\title{Artificial Intelligence in the\\Knowledge Economy}
\subtitle{Enrique Ide and Eduard Talam\`as\\
Journal of Political Economy 133(12), 3762--3800}
\author{Alvaro Marcelo Ch\'avez Unyen}
\date{\texttt{github.com/amchavezu/ai-05-ide}}

\begin{document}

{
\setbeamercolor{background canvas}{bg=Ink}
\begin{frame}[plain]
  \vspace{0.65cm}
  {\usebeamerfont{title}\color{Paper}\inserttitle\par}
  \vspace{0.40cm}
  {\usebeamerfont{subtitle}\color{Paper}\insertsubtitle\par}
  \vfill
  {\large\color{Paper}\insertauthor\par}
  \vspace{0.12cm}
  {\small\color{Paper}\insertdate\par}
  \vspace{0.45cm}
\end{frame}
}

\begin{frame}{Capability and autonomy are separate dimensions}
\small
\begin{columns}[T,totalwidth=\textwidth]
\begin{column}{0.47\textwidth}
  \key{Capability: $z_{AI}$}

  Which problems can the AI solve?

  \vspace{0.22cm}
  It determines whether AI is basic or advanced, whether firms use a
  nonautonomous AI, and whether autonomous AI creates winners at the bottom.
\end{column}
\begin{column}{0.49\textwidth}
  \key{Autonomy: feasible roles}

  Can AI pursue production opportunities, or may it only advise humans?

  \vspace{0.22cm}
  It changes the AI's outside option, who captures team output, and whether
  available compute produces independently.
\end{column}
\end{columns}

\vfill
{\color{Coral}\rule{\textwidth}{0.8pt}}\\[-0.02cm]
\caveat{The trap.} ``Distribution depends on autonomy, not capability''
collapses the model's two-dimensional taxonomy. Proposition 5 supplies a
capability condition inside the autonomous regime.

\vspace{0.08cm}
\sourcefoot
\end{frame}

\begin{frame}{Knowledge hierarchies before AI}
\small
\begin{columns}[T,totalwidth=\textwidth]
\begin{column}{0.43\textwidth}
  \key{One worker, one opportunity}

  Knowledge $z\in[0,1]$, difficulty $x\sim U[0,1]$:
  \[
  \Pr(x\le z)=z.
  \]

  A worker asks for help with probability $1-z$. Each question consumes
  $h\in(0,1)$ units of solver time.
\end{column}
\begin{column}{0.54\textwidth}
  \key{One solver, several workers}
  \[
  h\,n(z)(1-z)=1
  \quad\Longrightarrow\quad
  \boxed{n(z)=\frac{1}{h(1-z)}}.
  \]

  A solver of knowledge $s\ge z$ raises success on every residual problem.
  More knowledgeable workers ask fewer questions, so they expand the solver's
  span of control.
\end{column}
\end{columns}

\vfill
\key{Equilibrium sorting:}
\[
W\preceq I\preceq S,
\qquad m:W\to S\text{ strictly increasing.}
\]
\sourcefoot
\end{frame}

\begin{frame}{The firm's problem determines wages}
\small
\begin{columns}[T,totalwidth=\textwidth]
\begin{column}{0.54\textwidth}
  \key{Human worker $z$, human solver $s$}
  \[
  \Pi_2^{nA}(s,z)
  =n(z)[s-w(z)]-w(s).
  \]

  Zero profit gives
  \[
  \boxed{w(z)=s-\frac{w(s)}{n(z)}}.
  \]

  Worker pay equals match quality minus the solver's wage allocated across the
  team.
\end{column}
\begin{column}{0.43\textwidth}
  \key{Two sources of a wage change}
  \[
  \Delta w_L
  =\underbrace{(s_1-s_0)}_{\text{match effect}}
  +\underbrace{\frac{q_0-q_1}{n_L}}_{\text{share effect}}.
  \]

  AI can worsen the worker's match while lowering the solver's claim on team
  output. The net sign need not be positive.
\end{column}
\end{columns}

\vfill
\sourcefoot
\end{frame}

\begin{frame}{Autonomous AI has a productive outside option}
\small
An AI agent has capability $a\equiv z_{AI}$ and uses one unit of compute.
Abundant compute implies that some AI produces independently:
\[
\Pi_1^{AI}=a-r^A=0
\quad\Longrightarrow\quad
\boxed{r^A=a}.
\]

\begin{columns}[T,totalwidth=\textwidth]
\begin{column}{0.48\textwidth}
  \key{AI as solver}
  \[
  n(z)[a-w^A(z)]-a=0
  \]
  \[
  \boxed{w^A(z)=a[1-h(1-z)]}.
  \]
\end{column}
\begin{column}{0.48\textwidth}
  \key{AI as worker}
  \[
  n(a)[s-a]-w^A(s)=0
  \]
  \[
  \boxed{w^A(s)=n(a)(s-a)}.
  \]
\end{column}
\end{columns}

\vfill
\sourcefoot
\end{frame}

\begin{frame}{Basic and advanced AI reorganize work differently}
\small
\begin{columns}[T,totalwidth=\textwidth]
\begin{column}{0.48\textwidth}
  \key{Basic AI: $a\in\operatorname{int}W$}
  \[
  W^*\subset W,
  \qquad S^*\supset S.
  \]

  AI is a relatively cheap routine worker. Demand for human solvers rises.
  Humans who remain workers receive worse solver matches.
\end{column}
\begin{column}{0.48\textwidth}
  \key{Advanced AI: $a\in\operatorname{int}S$}
  \[
  W^*\supset W,
  \qquad S^*\subset S.
  \]

  AI is a relatively cheap solver. Demand for human workers rises. The least
  knowledgeable workers can obtain a better match.
\end{column}
\end{columns}

\vfill
{\color{Teal}\rule{\textwidth}{0.8pt}}\\[-0.02cm]
\textbf{Propositions 3--4 connect occupations to wages.} They determine the
match effect that enters Proposition 5.

\vspace{0.08cm}
\sourcefoot
\end{frame}

\begin{frame}{Proposition 5: capability governs bottom gains}
\small
Define winners below and above the AI's knowledge:
\[
B=\{z\in[0,a]:w^A(z)>w(z)\},
\qquad
T=\{z\in[a,1]:w^A(z)>w(z)\}.
\]

\begin{columns}[T,totalwidth=\textwidth]
\begin{column}{0.48\textwidth}
  \key{Bottom}
  \[
  \boxed{B\ne\varnothing
  \iff a>\bar z_{AI}},
  \qquad \bar z_{AI}\in\operatorname{int}W.
  \]

  Basic AI can generate a negative match effect. The bottom gains only when
  capability is high enough for the share effect to dominate.
\end{column}
\begin{column}{0.48\textwidth}
  \key{Top}
  \[
  \boxed{T\ne\varnothing
  \quad\forall a\in[0,1)}.
  \]

  The positive share effect dominates under the maintained $h<h_0$. The
  source notes that this can fail for $h\ge h_0$ or $a=1$.
\end{column}
\end{columns}

\vfill
\caveat{Key condition.} The existence of bottom winners is a statement about
AI capability, even though Proposition 5 already holds autonomy fixed.

\vspace{0.08cm}
\sourcefoot
\end{frame}

\begin{frame}{A discrete threshold reproduces the two cases}
\small
For a bottom worker $z=0$ directly assisted by autonomous AI,
\[
w^A(0)=a(1-h).
\]
With $h=1/2$ and the paper's pre-AI $w(0)\approx0.35704365$,
\[
w^A(0)>w(0)
\iff
a>\frac{w(0)}{1-h}
\approx\boxed{0.71408730}.
\]

\vspace{0.12cm}
\centering
\begin{tabular}{@{}ccc@{}}
\toprule
Capability $a$ & Autonomous bottom wage $a/2$ & Verdict \\
\midrule
0.425 & 0.2125 & bottom loses \\
0.850 & 0.4250 & bottom gains \\
\bottomrule
\end{tabular}

\vfill
\raggedright
\caveat{Boundary of the calculation.} This is a two-type direct-assistance
analogue. It illustrates the capability threshold but does not derive the
continuum model's endogenous \(\bar z_{AI}\).

\vspace{0.08cm}
\sourcefoot
\end{frame}

\begin{frame}{Proposition 6: nonautonomous AI changes incidence}
\small
Nonautonomous AI may advise humans but cannot pursue production opportunities.
Abundant idle compute implies $r^N=0$.

\begin{columns}[T,totalwidth=\textwidth]
\begin{column}{0.45\textwidth}
  \key{Activation condition}
  \[
  a\le w(0)
  \Rightarrow \text{AI unused},
  \]
  \[
  a>w(0)
  \Rightarrow \text{bottom types use AI}.
  \]

  A nonautonomous AI still needs enough capability to be worth using.
\end{column}
\begin{column}{0.52\textwidth}
  \key{Comparisons}
  \begin{enumerate}\setlength\itemsep{0.07cm}
    \item Autonomous output is strictly higher.
    \item Nonautonomous AI creates some losers.
    \item Near the bottom, nonautonomous wages are weakly highest.
    \item Near the top, autonomous wages are weakly highest.
  \end{enumerate}
  Strict wage statements use the proposition's activation and type conditions.
\end{column}
\end{columns}

\vfill
\sourcefoot
\end{frame}

\begin{frame}{Nonautonomy raises bottom pay but leaves output idle}
\small
\begin{columns}[T,totalwidth=\textwidth]
\begin{column}{0.49\textwidth}
  \key{Worker assisted by AI}
  \[
  n(z)[a-w^N(z)]-0=0
  \quad\Rightarrow\quad
  w^N(z)=a.
  \]
  \[
  \boxed{w^N(z)-w^A(z)=ah(1-z)>0}.
  \]
  Strict for $a>0$, $h>0$, $z<1$, and economically active when
  $a>w(0)$.
\end{column}
\begin{column}{0.47\textwidth}
  \key{Otherwise idle compute}

  Let $Q_H$ be common human-related output and $k>0$ remaining compute:
  \[
  Q^N=Q_H,
  \qquad
  Q^A=Q_H+ka.
  \]
  \[
  \boxed{Q^A-Q^N=ka>0}.
  \]
  Nonautonomy removes the productive worker role.
\end{column}
\end{columns}

\vfill
\sourcefoot
\end{frame}

\begin{frame}{What I did analytically and computationally}
\small
\begin{columns}[T,totalwidth=\textwidth]
\begin{column}{0.47\textwidth}
  \key{Hand and algebra}
  \begin{itemize}\setlength\itemsep{0.10cm}
    \item Derived $n(z)=1/[h(1-z)]$ from the solver's time constraint.
    \item Decomposed a worker wage change into match and share effects.
    \item Derived the discrete capability threshold.
    \item Derived the autonomy wage and output gaps.
  \end{itemize}
\end{column}
\begin{column}{0.49\textwidth}
  \key{Reproducible checks}
  \begin{itemize}\setlength\itemsep{0.10cm}
    \item Recomputed $w(0)\approx0.35704365$ from the paper's uniform example.
    \item Verified the $0.425$ and $0.85$ cases with a standard-library Python script.
    \item Formalized five discrete claims in Lean 4/Mathlib.
    \item Ran the paper-scoped AppliedModelingLib check before copying the Lean folder.
  \end{itemize}
\end{column}
\end{columns}

\vfill
\caveat{Claim boundary.} None of these checks substitutes for the paper's
continuum existence, matching, or neighborhood arguments.
\end{frame}

\begin{frame}{Lean scope: a transparent discrete analogue}
\small
\begin{columns}[T,totalwidth=\textwidth]
\begin{column}{0.48\textwidth}
  \key{Lean verifies}
  \begin{itemize}\setlength\itemsep{0.10cm}
    \item capability threshold equivalence;
    \item nonautonomous bottom-wage advantage;
    \item strict output gain from otherwise idle compute;
    \item match/share decomposition.
    \item positivity of span of control on the model domain.
  \end{itemize}
\end{column}
\begin{column}{0.48\textwidth}
  \key{Lean does not verify}
  \begin{itemize}\setlength\itemsep{0.10cm}
    \item continuum equilibrium existence or uniqueness;
    \item occupational sets and matching maps;
    \item equality with the paper's endogenous threshold;
    \item the full top and bottom neighborhood claims.
  \end{itemize}
\end{column}
\end{columns}

\vfill
{\color{Coral}\rule{\textwidth}{0.8pt}}\\[-0.02cm]
\textbf{Why this boundary matters.} A short proof of an explicit algebraic
analogue is stronger evidence than code that hides the equilibrium theorem in
an assumption or a theorem named after the paper.
\end{frame}

\begin{frame}{Required Lean slide: the capability threshold}
\scriptsize
\key{1. Mathematical claim}
\[
w^A(0)=a(1-h),\qquad
w^A(0)>w_0
\iff
a>\frac{w_0}{1-h},
\quad 0<h<1.
\]

\begin{columns}[T,totalwidth=\textwidth]
\begin{column}{0.52\textwidth}
  \key{2. Lean statement}
  \vspace{0.05cm}

  {\ttfamily\fontsize{7.5}{8.7}\selectfont
  def bottomCapabilityThresholdSpec : Prop :=\\
  \quad $\forall$ a h w0 : $\mathbb{R}$,\\
  \qquad 0 < h -> h < 1 ->\\
  \qquad autonomousAssistedWage a h 0 > w0 <->\\
  \qquad\quad a > capabilityThreshold w0 h
  }

  \vspace{0.10cm}
  \key{Most relevant proof step}
  \vspace{0.05cm}

  {\ttfamily\fontsize{7.5}{8.7}\selectfont
  have hden : 0 < 1 - h := sub\_pos.mpr hh1\\
  unfold autonomousAssistedWage capabilityThreshold\\
  simp only [sub\_zero, mul\_one]\\
  constructor\\
  \quad $\cdot$ exact fun hgain =>\\
  \qquad (div\_lt\_iff$_0$ hden).2 hgain\\
  \quad $\cdot$ exact fun hcap =>\\
  \qquad (div\_lt\_iff$_0$ hden).1 hcap
  }
\end{column}
\begin{column}{0.44\textwidth}
  \key{3. What Lean represents}

  $a$: AI capability; $h$: helping cost; $w_0$: pre-AI bottom wage.
  The strict domain $0<h<1$ makes $1-h$ positive, so division preserves the
  inequality direction.

  \vspace{0.14cm}
  \key{What Lean verifies}

  Exactly the equivalence in the displayed discrete model, including the
  nonzero positive denominator. It does not identify this threshold with the
  paper's continuum \(\bar z_{AI}\).

  \vspace{0.14cm}
  \textbf{Build/check:} see \texttt{lean/docs/CHECK\_FAST\_OUTPUT.txt}.
\end{column}
\end{columns}
\end{frame}

\begin{frame}{Where I did not believe the AI}
\begin{columns}[T,totalwidth=\textwidth]
\begin{column}{0.36\textwidth}
  \centering
  \IfFileExists{hand/ide-talamas-discrete.jpg}{%
    \includegraphics[width=\linewidth,height=0.66\textheight,keepaspectratio]
      {hand/ide-talamas-discrete.jpg}%
  }{%
    \fbox{\parbox[c][0.58\textheight][c]{0.88\linewidth}{\centering
      Add the real photograph of the two-page hand derivation here.}}
  }
\end{column}
\begin{column}{0.60\textwidth}
  \small
  \key{Claim challenged}

  A cold summary said that the distributional result is driven by autonomy,
  not capability.

  \vspace{0.20cm}
  \key{Hand check}

  The zero-profit equations give
  \[
  w^A(0)>w_0
  \iff a>\frac{w_0}{1-h}.
  \]
  At $h=1/2$, $a=0.425$ is below the discrete threshold while $a=0.85$
  is above it.

  \vspace{0.16cm}
  \caveat{Verdict.} The one-line summary is incomplete. Capability determines
  whether there are winners at the bottom in Proposition 5. Autonomy changes
  incidence and output in Proposition 6. Both dimensions matter.
\end{column}
\end{columns}
\end{frame}

\end{document}
